\section{Conclusion}
\label{sec:conclusion}

The \ar{} plan occupies a legally optimal position in the compactness
distribution.  It is more compact than the average valid plan ($61$st--$72$nd
percentile), confirming that compactness is a genuine output of the algorithm.
It is not at the extreme of the distribution (below the 95th percentile in
all states), avoiding the packing accusation that attaches to maximally
compact plans in sorted states.  And its compactness is traceable to a
stated, publicly documented objective (METIS edge-cut minimisation) rather
than post-hoc selection.

The compactness-partisan correlation within the \recom{} ensemble
($\rho \approx -0.1$ to $-0.3$ in sorted states) confirms the Rodden
effect: more compact plans tend to produce fewer Democratic seats in sorted
states.  This correlation is a property of geography, not of the \ar{}
algorithm.  Any compact algorithm applied to a sorted state would produce
a similar correlation.

The legal implication: the \ar{} plan's compactness is not an indicator of
gerrymandering.  Rather, it is evidence of the opposite.  A plan that is
more compact than average but not at the extreme, produced by an algorithm
with a public objective function and no partisan inputs, is precisely what
redistricting reform advocates have sought.  The ensemble comparison confirms
that this plan is not an outlier---it is a typical compact plan, made unusual
only by its verifiable, deterministic provenance.

Together with G.1 and G.2, this paper completes the three-dimensional
characterisation of the \ar{} plan's position in the ensemble:
compact but not an outlier, partisan-neutral in competitive states,
and consistent with minority representation requirements.  This
three-dimensional certificate constitutes the strongest available quantitative
evidence that the \ar{} plan is geographically reasonable and legally
defensible.
